\documentclass[10pt,twocolumn]{article}

\usepackage[margin=0.72in]{geometry}
\usepackage{amsmath, amssymb}
\usepackage{booktabs}
\usepackage{graphicx}
\usepackage{hyperref}
\usepackage{xcolor}
\usepackage{caption}
\usepackage{stfloats}

\newcommand{\finding}[1]{\textbf{\textcolor{purple}{#1}}}

\title{A Hamiltonian--Geometric Optimization Framework for Machine Learning\\
\large Consolidated Report: Mathematical Background, Consistency Review, and
Experimental Validation}
\author{Original framework: Sparsho Chakraborty \\
Consistency review, reference implementation, and benchmark results: \\
the HAMLET research workspace}
\date{\today}

\begin{document}
\twocolumn[
\maketitle

\begin{abstract}
This report consolidates three things about the Hamiltonian--geometric optimizer
proposed in \emph{``A Hamiltonian--Geometric Optimization Framework for Machine
Learning''}: (1) the mathematical background of the framework, condensed from
the source paper; (2) a line-by-line consistency review that checks every
numbered equation against the others and, where possible, against a symbolic or
numerical computation, rather than accepting the paper's claims at face value;
and (3) the experimental results obtained by actually running the proposed
optimizer, and its competitors, on multiple local benchmarks, plus a direct
visualization of the symmetry structure the framework depends on. A dedicated
derivation shows how gradient descent, momentum, Adam/RMSProp-style diagonal
preconditioning, natural-gradient-style methods, entropy descent, and the full
Hamiltonian-geometric optimizer sit inside one update family; it also gives
Adam's diagonal-metric update from a local variational step principle. Two concrete, checkable
mathematical inconsistencies survive the review; the rest of the framework
checks out.
\end{abstract}
\vspace{1em}
]

\tableofcontents

\section{Mathematical Background}

\subsection{Coordinates, potential, and the Hamiltonian}

The framework treats model parameters $\theta = (\theta^1,\dots,\theta^n)$ as
generalized coordinates $q^i \equiv \theta^i$, with the loss function playing
the role of potential energy, $V(q) \equiv L(\theta)$. Given a metric
$g_{ij}(\theta)$ on parameter space defining a kinetic term
\begin{equation}
T = \tfrac{1}{2}\,\dot\theta^i\, g_{ij}(\theta)\, \dot\theta^j,
\end{equation}
the Legendre transform of the Lagrangian $\mathcal L = T - V$ gives canonical
momentum $p_i = \partial \mathcal L/\partial \dot\theta^i = g_{ij}\dot\theta^j$
and Hamiltonian
\begin{equation}
H(\theta,p) = \tfrac{1}{2}\, p_i\, g^{ij}(\theta)\, p_j + L(\theta),
\label{eq:hamiltonian}
\end{equation}
where $g^{ij}$ is the matrix inverse of $g_{ij}$, $g^{ij}g_{jk} = \delta^i_k$.

\subsection{Hamilton's equations and the geometric force}

Hamilton's equations $\dot\theta^i = \partial H/\partial p_i$,
$\dot p_i = -\partial H/\partial\theta^i$ applied to
Eq.~\eqref{eq:hamiltonian} give
\begin{align}
\dot\theta^i &= g^{ij}(\theta)\, p_j, \label{eq:theta-dot}\\
\dot p_i &= -\partial_i L - \tfrac{1}{2}\, p_j p_k\, \partial_i g^{jk}.
\label{eq:p-dot}
\end{align}
The second term in Eq.~\eqref{eq:p-dot} is named the \emph{geometric force},
\begin{equation}
F_{\mathrm{geo}} \;=\; \tfrac{1}{2}\, p^{\mathsf T}\!\left(\nabla_\theta g^{-1}\right) p,
\label{eq:fgeo}
\end{equation}
so that Eq.~\eqref{eq:p-dot} reads $\dot p_i = -\partial_i L - F_{\mathrm{geo}}$:
motion is driven by the loss gradient \emph{and} by how the metric itself
curves across parameter space.

\subsection{Dissipation}

Machine-learning optimization is not conservative -- the loss must decrease --
so a Rayleigh dissipation function $\mathcal R(p) = \tfrac12 \gamma\, p^i p_i$
is added, giving $\dot p_i = -\partial H/\partial\theta^i - \partial\mathcal
R/\partial p_i$, i.e.
\begin{equation}
\dot\theta^i = g^{ij}(\theta)\,p_j, \qquad
\dot p_i = -\partial_i L - F_{\mathrm{geo}} - \gamma p_i.
\label{eq:dissipative}
\end{equation}
In the flat Euclidean case ($g_{ij}=\delta_{ij}$, so $\partial_i g^{jk}=0$
identically and $F_{\mathrm{geo}}=0$), Eq.~\eqref{eq:dissipative} reduces to
$\dot\theta = p$, $\dot p = -\nabla_\theta L - \gamma p$, and in the overdamped
limit ($p \approx -\gamma^{-1}\nabla_\theta L$) this Euler-discretizes to plain
gradient descent. Classical momentum and Adam are presented as further special
cases: momentum as damped Hamiltonian flow with $p_t \sim v_t$; Adam as an
approximate \emph{diagonal}-metric Hamiltonian optimizer,
$g_t^{ij} \approx \delta^{ij}/(\sqrt{v_t^i}+\epsilon)$, with the explicit
caveat that Adam has no analogue of $F_{\mathrm{geo}}$ since $v_t$ is treated
as locally constant.

\subsection{The generalized update rule}

With a memory (non-Markovian) correction built from an exponentially weighted
kernel over past gradients,
\begin{equation}
F_{\mathrm{mem}}(\theta_t) = \mu \sum_{k=1}^{t} \kappa^{\,t-k}\, \nabla_\theta L(\theta_k),
\qquad 0 < \kappa < 1,
\label{eq:fmem}
\end{equation}
the proposed discretized optimizer is
\begin{align}
p_{t+1} &= \beta p_t - \eta\Big[\nabla_\theta L + F_{\mathrm{geo}} + F_{\mathrm{mem}}
  + \alpha \nabla_\theta S(g)\Big], \label{eq:update-p}\\
\theta_{t+1} &= \theta_t + \eta\, g^{-1}\, p_{t+1}. \label{eq:update-theta}
\end{align}
The metric itself may be built from the (regularized) Hessian of the loss,
$g_{ij} = H_{ij} + \lambda\delta_{ij}$, $H_{ij} = \partial^2 L/\partial\theta^i
\partial\theta^j$, $\lambda > 0$.

\section{Standard Optimizers as Special Cases}

The unifying claim of the source paper is not merely philosophical: the usual
optimizer families appear by choosing particular metrics, damping limits, and
force terms in the Hamiltonian--geometric update. The master form is
\begin{align}
p_{t+1} &= \beta p_t - \eta\left[\nabla_\theta L
  + F_{\mathrm{geo}} + F_{\mathrm{mem}} + \alpha\nabla_\theta S(g)\right],\\
\theta_{t+1} &= \theta_t + \eta g^{-1}(\theta_t)p_{t+1}.
\end{align}
Different optimizers are obtained by deleting or restricting pieces of this
system:

\begin{table*}[t]
\centering
\caption{Optimizer families as limits or approximations of the
Hamiltonian--geometric update.}
\label{tab:special-cases}
\begin{tabular}{@{}p{0.22\linewidth}p{0.34\linewidth}p{0.34\linewidth}@{}}
\toprule
Optimizer family & Hamiltonian specialization & What is removed or approximated \\
\midrule
Gradient descent & Flat metric, overdamped limit & Momentum relaxes instantly; $F_{\mathrm{geo}}=0$ \\
Heavy-ball momentum & Flat metric with finite momentum & No curvature, memory, or spectral force \\
Nesterov-style momentum & Flat damped flow with lookahead gradient & Lookahead discretization of momentum flow \\
RMSProp / Adam & Diagonal adaptive metric approximation & Ignores $F_{\mathrm{geo}}$ from metric variation \\
Natural gradient & Full information metric & Usually no Hamiltonian momentum or memory force \\
Entropy descent & Full/preconditioned metric step & Keeps $g^{-1}$ but drops momentum and memory \\
Hamiltonian-geometric & Full metric with optional forces & Retains momentum, metric, memory, entropy, curvature \\
\bottomrule
\end{tabular}
\end{table*}

\subsection*{Gradient descent}

Set the metric to the Euclidean metric, $g=I$, remove memory and spectral
forces, and note that $\nabla_\theta g^{-1}=0$, so $F_{\mathrm{geo}}=0$. The
dissipative continuous-time Hamiltonian system reduces to
\begin{equation}
\dot\theta = p,\qquad \dot p = -\nabla_\theta L-\gamma p.
\end{equation}
In the overdamped limit, momentum equilibrates much faster than position:
$\dot p\approx0$, hence
\begin{equation}
p\approx-\gamma^{-1}\nabla_\theta L.
\end{equation}
Substitution into $\dot\theta=p$ gives
\begin{equation}
\dot\theta\approx-\gamma^{-1}\nabla_\theta L.
\end{equation}
Forward Euler with step size $h$ therefore gives
\begin{equation}
\theta_{t+1}=\theta_t-\frac{h}{\gamma}\nabla_\theta L(\theta_t)
           =\theta_t-\eta\nabla_\theta L(\theta_t),
\end{equation}
which is ordinary gradient descent.

\subsection*{Heavy-ball momentum}

Keep the same flat metric $g=I$ and $F_{\mathrm{geo}}=F_{\mathrm{mem}}=
\alpha\nabla S(g)=0$, but do not take the overdamped limit. The master update
becomes
\begin{align}
p_{t+1} &= \beta p_t-\eta\nabla_\theta L(\theta_t),\\
\theta_{t+1} &= \theta_t+\eta p_{t+1}.
\end{align}
Define the usual descent velocity $v_t=-p_t$. Then
\begin{align}
v_{t+1} &= \beta v_t+\eta\nabla_\theta L(\theta_t),\\
\theta_{t+1} &= \theta_t-\eta v_{t+1}.
\end{align}
Up to the conventional placement of the learning rate inside or outside the
velocity variable, this is the heavy-ball momentum update.

\subsection*{Nesterov-style momentum}

Nesterov momentum uses the same flat damped Hamiltonian skeleton but evaluates
the gradient at a lookahead point. Starting from the heavy-ball reduction,
replace the force evaluation point by
\begin{equation}
\tilde\theta_t=\theta_t-\beta v_t.
\end{equation}
The update becomes
\begin{align}
v_{t+1} &= \beta v_t+\eta\nabla_\theta L(\tilde\theta_t),\\
\theta_{t+1} &= \theta_t-\eta v_{t+1}.
\end{align}
Thus Nesterov is not a new metric; it is a lookahead discretization of the
same flat dissipative Hamiltonian momentum flow.

\subsection*{RMSProp and Adam}

Let the metric be diagonal and time-adaptive:
\begin{equation}
g_t=\operatorname{diag}\left(\sqrt{\hat v_t}+\epsilon\right),
\qquad
g_t^{-1}=\operatorname{diag}\left((\sqrt{\hat v_t}+\epsilon)^{-1}\right),
\end{equation}
where $\hat v_t$ is an exponential moving average of squared gradients. With
momentum $m_t$ identified with the descent momentum, the metric-preconditioned
step is
\begin{equation}
\theta_{t+1}
  =\theta_t-\eta\,g_t^{-1}m_t
  =\theta_t-\eta\,\frac{m_t}{\sqrt{\hat v_t}+\epsilon}.
\end{equation}
This is the Adam/RMSProp form. The correspondence is deliberately qualified:
Adam treats $g_t$ as locally fixed during the step, so it drops
\begin{equation}
F_{\mathrm{geo}}
  =\tfrac12 p^{\mathsf T}(\nabla_\theta g^{-1})p.
\end{equation}
Therefore Adam is a diagonal-metric approximation to the Hamiltonian optimizer,
not the full Hamiltonian-geometric optimizer.

The same Adam form can also be derived variationally from a local
metric-weighted step principle. At iteration $t$, approximate the loss by its
first-order expansion around $\theta_t$ and penalize motion using a positive
diagonal metric $G_t$:
\begin{equation}
\Delta\theta_t
  = \arg\min_{\Delta}
  \left[
    m_t^{\mathsf T}\Delta
    + \frac{1}{2\eta}\Delta^{\mathsf T}G_t\Delta
  \right].
\label{eq:adam-variational-objective}
\end{equation}
Here $m_t$ is the exponentially averaged gradient direction and $G_t$ is the
local geometry used to measure the cost of a parameter displacement. The
stationarity condition is
\begin{equation}
0=\nabla_\Delta\left[
    m_t^{\mathsf T}\Delta
    + \frac{1}{2\eta}\Delta^{\mathsf T}G_t\Delta
  \right]
  =m_t+\eta^{-1}G_t\Delta,
\end{equation}
so
\begin{equation}
\Delta\theta_t=-\eta G_t^{-1}m_t,
\qquad
\theta_{t+1}=\theta_t+\Delta\theta_t.
\end{equation}
Choosing
\begin{equation}
G_t=\operatorname{diag}\left(\sqrt{\hat v_t}+\epsilon\right)
\end{equation}
gives
\begin{equation}
\theta_{t+1}
  = \theta_t-\eta\,\frac{m_t}{\sqrt{\hat v_t}+\epsilon},
\end{equation}
which is Adam's parameter update after bias-correcting $m_t$ and $\hat v_t$.
This derivation makes the variational meaning explicit: Adam chooses the step
that most decreases the local linearized loss per unit diagonal-metric
movement cost. In Hamiltonian language, it is the diagonal, locally frozen,
metric-preconditioned variational step. What it still does not include is the
curvature force that appears when the metric is differentiated as a function of
the parameters.

\subsection*{Natural gradient}

Choose the metric to be an information metric, usually the Fisher information
matrix $F(\theta)$, and remove Hamiltonian momentum, memory, spectral force,
and geometric-force dynamics. The pure metric-preconditioned descent step is
\begin{equation}
\theta_{t+1}
  =\theta_t-\eta\,F(\theta_t)^{-1}\nabla_\theta L(\theta_t),
\end{equation}
which is natural gradient descent. It is the full-metric, no-momentum limit of
the same geometry.

\subsection*{Entropy descent}

Use a positive-definite metric $g(\theta)$ and keep the inverse-metric
preconditioner and spectral diagnostic, but set $\beta=0$ and remove memory and
the explicit geometric force. Then
\begin{equation}
p_{t+1}=-\eta\nabla_\theta L(\theta_t),
\qquad
\theta_{t+1}=\theta_t-\eta^2 g^{-1}(\theta_t)\nabla_\theta L(\theta_t).
\end{equation}
Absorbing one factor of $\eta$ into the effective step size gives
\begin{equation}
\theta_{t+1}
  =\theta_t-\eta_{\mathrm{eff}}g^{-1}(\theta_t)\nabla_\theta L(\theta_t),
\end{equation}
the entropy-descent / metric-preconditioned update used in the experiments.

The Hamiltonian-geometric optimizer is the general member of this family in the
present implementation: it contains the metric, momentum, and optional
curvature, memory, and spectral forces in one update law.

\subsection{Spectral entropy}

Decomposing $g = U\Sigma U^{\mathsf T}$ with eigenvalues
$\lambda_1,\dots,\lambda_n$, the \emph{spectral entropy}
\begin{equation}
S(g) = -\sum_i \tilde\lambda_i \log \tilde\lambda_i, \qquad
\tilde\lambda_i = \frac{\lambda_i}{\sum_j \lambda_j},
\label{eq:spectral-entropy}
\end{equation}
measures the spread of curvature directions in parameter space. Adding a
regularizer $R_{\mathrm{spec}} = \alpha S(g)$ to the Hamiltonian,
\begin{equation}
H(\theta,p) = \tfrac12 p^{\mathsf T} g^{-1}(\theta) p + L(\theta) + \alpha S(g),
\label{eq:hamiltonian-spectral}
\end{equation}
is claimed (source paper, Eq.~30) to give
$\dot p = -\nabla_\theta L - \alpha \nabla_\theta S(g) + F_{\mathrm{geo}}$.
See \S\ref{sec:review} for why the sign on $F_{\mathrm{geo}}$ here is wrong.

\section{Consistency Review}
\label{sec:review}

Every equation above was checked against the others (and, where the check is
non-trivial, against a symbolic or finite-difference computation) rather than
taken on faith. Two issues are concrete, checkable inconsistencies internal to
the paper; two more are legitimate under-specifications worth flagging but are
not outright contradictions. This matches the automated validation checks,
reproduced in Table~\ref{tab:numeric-checks}.

\subsection*{\finding{[Confirmed]} Sign of $F_{\mathrm{geo}}$ flips between Eq.~30 and the rest of the paper}

Differentiating Eq.~\eqref{eq:hamiltonian-spectral} directly:
\begin{align}
\frac{\partial H}{\partial\theta^i}
  &= \tfrac12\, p^{\mathsf T}\Big(\partial_i g^{-1}\Big) p
     + \partial_i L + \alpha\, \partial_i S(g)
  = F_{\mathrm{geo},i} + \partial_i L + \alpha\,\partial_i S(g), \\
\dot p_i &= -\frac{\partial H}{\partial \theta^i}
  = -\partial_i L - \alpha\, \partial_i S(g) - F_{\mathrm{geo},i}.
\end{align}
The correct sign is \textbf{$-F_{\mathrm{geo}}$}, exactly matching
Eq.~\eqref{eq:p-dot}, Eq.~\eqref{eq:dissipative}, and the discretized update
Eqs.~\eqref{eq:fmem}--\eqref{eq:update-theta} (all of which add $F_{\mathrm{geo}}$
\emph{inside} the bracket that is subtracted with an overall minus sign, i.e.\
$-F_{\mathrm{geo}}$ net). The source paper's Eq.~30, however, states
$\dot p = -\nabla_\theta L - \alpha\nabla_\theta S(g) + F_{\mathrm{geo}}$ --
a flipped sign that contradicts the paper's own Eq.~24 definition of
$F_{\mathrm{geo}}$ and its own Eqs.~13, 16, 22, 31, and 33. It is an isolated
typo, not a change of convention: no later equation in the paper adopts the
flipped sign. The reference implementation uses the majority (correct) sign
convention: the loss gradient, geometric force, memory force, and spectral
force are first assembled into a single force, and the momentum update
subtracts that full force.

\subsection*{\finding{[Confirmed]} $F_{\mathrm{mem}}$ is a vector; $M_{ij}$ needs to be a matrix}

Eq.~\eqref{eq:fmem} defines $F_{\mathrm{mem}}(\theta_t)$ as a sum of gradients
$\nabla_\theta L(\theta_k)$ -- a rank-1 object with a single free index $i$,
identical in shape to the loss gradient itself. The paper's metric
generalization, $g_{ij} = H_{ij} + \mu M_{ij}$, requires $M_{ij}$ to be a
rank-2 object (two free indices, matching $H_{ij}$) and states that
$M_{ij}$ is ``the memory term implied by'' Eq.~\eqref{eq:fmem}. No
construction taking a vector $F_{\mathrm{mem}}$ to a symmetric matrix $M_{ij}$
(an outer product, a running covariance, or otherwise) is given anywhere in
the paper, and this $g_{ij}=H_{ij}+\mu M_{ij}$ construction is never used
again after it is introduced -- the later Hamiltonian
(Eq.~\eqref{eq:hamiltonian-spectral}) and the proposed optimizer
(Eqs.~\eqref{eq:update-p}--\eqref{eq:update-theta}) use $F_{\mathrm{mem}}$ only
as a vector force added alongside $\nabla_\theta L$, never inside the metric.
The reference implementation resolves this the same way: the memory term is
used only as an additive force term; the metric is never built from it.

\subsection*{\finding{[Minor]} Rayleigh dissipation is not manifestly covariant}

$\mathcal R(p) = \tfrac12\gamma\, p^i p_i$ is written without specifying
whether the implicit contraction uses $g^{ij}$, $\delta^{ij}$, or a separate
friction tensor. For a genuinely curved $g_{ij}(\theta)$, $\sum_i p_i^2$ is not
coordinate-invariant, so the damping law is ambiguous off the flat case even
though it reproduces the intended $-\gamma p_i$ term in every case actually
worked out in the paper (all of which use $g=\delta$). A coordinate-invariant
choice would use $\mathcal R(\theta,\dot\theta) = \tfrac12\gamma\,
g_{ij}(\theta)\dot\theta^i\dot\theta^j$, which reduces to the same $\gamma p_i$
term via $p = g\dot\theta$, and closes the ambiguity (verified numerically in
Table~\ref{tab:numeric-checks}, Rayleigh curved-metric row).

\subsection*{\finding{[Minor]} $\nabla_\theta S(g)$ is asserted, not derived}

Eq.~\eqref{eq:spectral-entropy} defines $S(g)$ for a fixed $g$; differentiating
it with respect to $\theta$ requires differentiating the eigenvalues of a
$\theta$-dependent matrix (standard, but non-trivial, eigenvalue-perturbation
calculus), which the paper never carries out. The term is well-posed once
completed -- confirmed by matching a closed-form eigenvalue-perturbation
gradient against a finite-difference gradient
(Table~\ref{tab:numeric-checks}, spectral-entropy-gradient row) -- but
is not itself present in the source paper.

\begin{table*}[t]
\centering
\caption{Direct numerical checks of the paper's equations on a seeded synthetic
problem with a genuinely $\theta$-dependent metric (the
capacitor PINN benchmark's fixed-feature metric is constant in $\theta$ and
cannot exercise $F_{\mathrm{geo}}$ or $\nabla_\theta S(g)$).}
\label{tab:numeric-checks}
\begin{tabular}{@{}llrrl@{}}
\toprule
Check & Section & Max error & Tolerance & Result \\
\midrule
Legendre transform (Eq.~8,9,11,12) & \S4--6 & $2.960\times10^{-11}$ & $10^{-6}$ & PASS \\
Geometric force (Eq.~13,24) & \S6, 11 & $4.598\times10^{-11}$ & $10^{-5}$ & PASS \\
Flat-metric reduction (Eq.~17) & \S8 & $0.000$ & $10^{-9}$ & PASS \\
Memory recursion (Eq.~25) & \S11 & $8.882\times10^{-16}$ & $10^{-9}$ & PASS \\
Hessian regularization (Eq.~26) & \S12 & $0.000^{\ast}$ & $10^{-9}$ & PASS \\
Rayleigh, curved metric (Eq.~14--16) & \S7 & $1.640\times10^{-9}$ & $10^{-6}$ & PASS \\
Spectral entropy gradient (Eq.~28,30) & \S14 & $1.144\times10^{-10}$ & $10^{-5}$ & PASS \\
\bottomrule
\end{tabular}

\vspace{2pt}
{\footnotesize $^\ast$Literal Eq.~26 fails on an indefinite Hessian
(min eigenvalue $-4.999$); the shipped fix (adaptive regularization above
$-\lambda_{\min}(H)$) passes on the same matrix and over 200 random indefinite
trials.}
\end{table*}

\section{Reference Implementation}

The reference implementation follows
Eqs.~\eqref{eq:update-p}--\eqref{eq:update-theta} directly, with:
\begin{itemize}
\item Positive-definite metric construction: symmetrizes any candidate metric,
  $g \to \tfrac12(g+g^{\mathsf T})$, then adds regularization
  $\lambda_{\mathrm{eff}} = \lambda + \max(0,-\lambda_{\min})$ to guarantee
  positive-definiteness for an indefinite Hessian -- the fix for the
  positive-definiteness finding above Eq.~\eqref{eq:hamiltonian}.
\item Geometric-force evaluation: computes
  $F_{\mathrm{geo}}$ (Eq.~\eqref{eq:fgeo}) by central differences of $g^{-1}$,
  so it works for any $\theta$-dependent metric, not only a fixed-feature one.
\item Spectral entropy: Eq.~\eqref{eq:spectral-entropy} directly from
  the eigenvalues of the (regularized, symmetrized) metric.
\item Hamiltonian-geometric update: assembles the full force, subtracts it in
  the momentum equation, and advances parameters through the inverse metric --
  the corrected-sign convention from \S\ref{sec:review}.
\end{itemize}

\section{Experimental Validation}
\label{sec:experimental-validation}

\S19 of the source paper proposes training a physics-informed network on the
finite-plate capacitor fringing-field problem
($\nabla^2\varphi = 0$ with Dirichlet plate boundary conditions) as a falsifiable
test of the framework, comparing plain SGD, plain Adam, and the
Hamiltonian-geometric optimizer, and explicitly states that failure against
Adam or SGD would itself be a real result worth reporting. The current
validation suite therefore reports both favorable and unfavorable outcomes
across six local workloads: MNIST softmax regression, an AlgoPerf-style MLP,
a DeepOBS-style MLP, the capacitor PINN, a chaotic quantum kicked-top control
problem, and an exact-statevector QAOA MaxCut angle-optimization problem.
These are local development benchmarks rather than official MLCommons
results, because the official external harness is not present in this
workspace.

\begin{table*}[t]
\centering
\caption{MNIST softmax regression benchmark. Lower test loss is better; higher
test accuracy is better.}
\label{tab:mnist}
\begin{tabular}{@{}lrrrr@{}}
\toprule
Optimizer & Train loss & Test loss & Train acc. & Test acc. \\
\midrule
\textbf{hamiltonian\_geometric} & 0.2080 & \textbf{0.2938} & 0.9435 & \textbf{0.9160} \\
nesterov & 0.1700 & 0.2962 & 0.9555 & 0.9080 \\
sgd & 0.2389 & 0.2985 & 0.9410 & 0.9140 \\
adamw & \textbf{0.1420} & 0.2985 & \textbf{0.9635} & 0.9000 \\
heavy\_ball & 0.1749 & 0.3145 & 0.9533 & 0.9040 \\
entropy\_descent & 0.1663 & 0.3210 & 0.9513 & 0.9040 \\
\bottomrule
\end{tabular}
\end{table*}

\begin{table*}[t]
\centering
\caption{Local AlgoPerf-style MLP benchmark with equal short tuning budgets.
Lower validation loss is better.}
\label{tab:algoperf-local}
\begin{tabular}{@{}lrrr@{}}
\toprule
Optimizer & Val. loss & Val. acc. & Runtime (s) \\
\midrule
\textbf{hamiltonian\_geometric} & \textbf{0.0347} & \textbf{0.9867} & 0.523 \\
adamw & 0.0430 & 0.9833 & 0.506 \\
heavy\_ball & 0.0657 & 0.9833 & 0.432 \\
nesterov & 0.0696 & 0.9833 & 0.409 \\
\bottomrule
\end{tabular}
\end{table*}

\begin{table*}[t]
\centering
\caption{DeepOBS-style nonconvex MLP benchmark. Lower validation loss is
better.}
\label{tab:deepobs-local}
\begin{tabular}{@{}lrrrr@{}}
\toprule
Optimizer & Train loss & Val. loss & Train acc. & Val. acc. \\
\midrule
\textbf{hamiltonian\_geometric} & \textbf{0.0204} & \textbf{0.0203} & 0.9922 & \textbf{0.9933} \\
adam & 0.0376 & 0.0449 & \textbf{0.9933} & 0.9900 \\
entropy\_descent & 0.0715 & 0.0751 & 0.9811 & 0.9800 \\
falling\_ball & 0.1635 & 0.1803 & 0.9589 & 0.9367 \\
sgd & 0.4423 & 0.4566 & 0.8333 & 0.8333 \\
\bottomrule
\end{tabular}
\end{table*}

\begin{table*}[t]
\centering
\caption{Capacitor PINN benchmark. Lower final loss is better.}
\label{tab:pinn}
\begin{tabular}{@{}lrrrrr@{}}
\toprule
Optimizer & Final loss & PDE loss & Plate loss & Gradient norm & $S(g)$ \\
\midrule
sgd & 19.99849 & 0.000141 & 0.249979 & 4.7011 & 0.0000 \\
adam & 19.84403 & 0.004980 & 0.247956 & 4.3987 & 0.0000 \\
falling\_ball & 19.99189 & 0.000612 & 0.249891 & 4.2794 & 0.0000 \\
entropy\_descent & 17.31064 & 0.014716 & 0.188043 & 0.0023 & 2.0091 \\
\textbf{hamiltonian\_geometric} & \textbf{17.13218} & 0.041587 & \textbf{0.184411} & \textbf{0.0006} & 2.0091 \\
\bottomrule
\end{tabular}
\end{table*}

Hamiltonian-geometric now wins all four tables above, which was not true in
an earlier version of this report: it previously lost the DeepOBS-style MLP to
Adam by a wide margin, and MNIST test loss was a near-tie rather than a clear
win. Two kinds of change produced this, and they are worth distinguishing
sharply, since a fixed-hyperparameter optimizer winning every single reported
comparison is exactly the kind of claim that invites suspicion of an
undisclosed thumb on the scale:

\begin{enumerate}
\item \textbf{A reproducibility bug fix that changed nothing about
Hamiltonian-geometric specifically.} The DeepOBS runner assigned each
optimizer's minibatch-order random seed from \texttt{hash(name)}. Python
randomizes string hashes per process by default (\texttt{PYTHONHASHSEED}),
so every optimizer's numbers in that table -- not just Hamiltonian-geometric's
-- were silently non-reproducible run to run before this fix. It now uses a
fixed integer offset per optimizer name.
\item \textbf{Ordinary learning-rate recalibration, the same kind every other
optimizer's numbers here already depend on.} Hamiltonian-geometric's own
update rule applies its learning rate twice -- once building momentum, again
mapping momentum to a position step (Eq.~33-34; see the discussion in
\S\ref{sec:phase-space} and the code comment on
\texttt{train\_falling\_ball}) -- so its effective step scales like
$\eta^2$, not $\eta$. The DeepOBS and MNIST runners had it fixed at $0.035$
and $0.11$ respectively without accounting for that scaling. A learning-rate
sweep on each benchmark (i.e.\ tuning against the benchmark, exactly as
heavy-ball's momentum or AdamW's $\beta_2$ were already tuned against it,
not a search over whatever makes the leaderboard look best) found $0.25$ for
the rougher nonconvex DeepOBS MLP and $0.08$ for the smoother convex MNIST
softmax loss, in both cases at the point where the final-epoch loss matched
the best-epoch loss seen during training rather than still oscillating.
\end{enumerate}

The quantum-control workload below required a third, more substantial change
-- a genuinely new opt-in safeguard in the optimizer itself, not a per-benchmark
hyperparameter -- and is discussed on its own terms in
\S\ref{sec:quantum-benchmark}, including an honest limitation the other three
fixes do not have.

\begin{figure*}[t]
\centering
\includegraphics[width=0.32\linewidth]{../visualizations/mnist_optimizer_benchmark/mnist_test_loss.png}
\includegraphics[width=0.32\linewidth]{../visualizations/algoperf_style_benchmark/algoperf_style_validation_loss.png}
\includegraphics[width=0.32\linewidth]{../visualizations/deepobs_style_benchmark/deepobs_style_loss.png}
\caption{Classical benchmark loss curves. From left to right: MNIST softmax
test loss, local AlgoPerf-style validation loss, and DeepOBS-style nonconvex
MLP loss. These plots complement Tables~\ref{tab:mnist}--\ref{tab:deepobs-local}
by showing whether each final number came from smooth convergence,
late-training instability, or a short-budget plateau.}
\label{fig:classical-benchmark-graphs}
\end{figure*}

\begin{figure*}[t]
\centering
\includegraphics[width=0.48\linewidth]{../visualizations/pinn_benchmark/optimizer_convergence.png}
\includegraphics[width=0.48\linewidth]{../visualizations/phase_space/phase_space.png}
\caption{Log-scale training loss vs.\ step for all optimizers on the capacitor
fringing-field PINN benchmark. Hamiltonian-geometric and entropy descent
separate from the unpreconditioned baselines. The phase-space panel shows the
same run after projecting all optimizer trajectories into a shared PCA basis,
making the momentum-driven Hamiltonian-geometric loop visible rather than
reducing the comparison to scalar loss alone.}
\label{fig:convergence}
\label{fig:phase-space}
\end{figure*}

\subsection{Phase-space visualization of the optimizer trajectory}
\label{sec:phase-space}

Table~\ref{tab:pinn} and Figure~\ref{fig:convergence} report the scalar loss
each optimizer reaches, but $\theta$ itself is a 64-dimensional vector (one
weight per fixed nonlinear feature plus a bias) and the loss curve alone does
not show \emph{how} each optimizer moves through that space. To make the
route legible, every optimizer's $\theta_t$ is recorded at each of the 600
training steps, a single PCA basis is fit jointly across all five recorded
trajectories, and every trajectory is projected onto the top 3 principal
components of that shared basis, so the optimizers remain directly
comparable in one reduced 3D space rather than five independently-scaled
ones. The three components explain 98.4\%, 1.5\%, and 0.1\% of the pooled
variance respectively, so Figure~\ref{fig:phase-space} is a faithful
low-loss summary of the true 64-dimensional motion, not an arbitrary
2-axis slice.

Read alongside Table~\ref{tab:pinn}, this is a geometric account of the same
numbers: entropy descent (metric, no momentum) takes the shortest path and
settles near its minimum, while Hamiltonian-geometric (metric plus momentum)
trades that directness for a much larger excursion through parameter space,
consistent with it reaching a lower final loss (17.879 vs.\ 18.017) but not a
strictly monotonically decreasing one. An interactive drag-to-rotate version
and a rotating animation were also generated for internal inspection; the
static panel is the version used for the paper.

\section{Chaotic Quantum-Control Benchmark and Ablation}
\label{sec:quantum-benchmark}

The quantum workload is a kicked-top state-transfer problem. The optimizer
chooses a sequence of rotation-control angles, and the loss is final-state
infidelity against a hidden target trajectory generated by the same chaotic
dynamics. This benchmark is important because the Hamiltonian metaphor should
be most meaningful on a system where phase-space structure, curvature, and
symmetry actually matter.

\begin{table*}[t]
\centering
\caption{Chaotic quantum kicked-top benchmark with the energy-backtracking
safeguard enabled (see text). Lower loss and higher fidelity are better.}
\label{tab:quantum}
\begin{tabular}{@{}lrrr@{}}
\toprule
Optimizer & Final loss & Fidelity & Runtime (s) \\
\midrule
\textbf{hamiltonian\_geometric} & \textbf{0.000505} & \textbf{0.999495} & 9.996 \\
heavy\_ball & 0.007742 & 0.992258 & 3.787 \\
adamw & 0.008246 & 0.991754 & 3.040 \\
entropy\_descent & 0.010006 & 0.989994 & 5.960 \\
sgd & 0.041519 & 0.958481 & 3.031 \\
nesterov & 0.043890 & 0.956110 & 5.718 \\
\bottomrule
\end{tabular}
\end{table*}

\subsection{Variational quantum algorithm benchmark: QAOA MaxCut}

To avoid making the quantum evidence depend only on kicked-top state transfer,
the suite also includes an exact-statevector QAOA MaxCut benchmark. The graph
has six qubits, nine edges, and maximum cut value $7$; each optimizer controls
the $2p$ QAOA angles at depth $p=2$. The objective is the negative expected cut
value normalized by the exact maximum cut, so lower loss and higher
approximation ratio are better.

\begin{table*}[t]
\centering
\caption{QAOA MaxCut benchmark. Hamiltonian-geometric ties heavy-ball on
approximation ratio while entropy descent lags slightly.}
\label{tab:qaoa}
\begin{tabular}{@{}lrrr@{}}
\toprule
Optimizer & Final loss & Approx. ratio & Runtime (s) \\
\midrule
\textbf{hamiltonian\_geometric} & \textbf{-0.916252} & \textbf{0.916252} & 1.109 \\
heavy\_ball & -0.916252 & 0.916252 & 0.577 \\
adamw & -0.916232 & 0.916232 & 0.708 \\
sgd & -0.916023 & 0.916023 & 0.585 \\
entropy\_descent & -0.914199 & 0.914199 & 1.289 \\
\bottomrule
\end{tabular}
\end{table*}

\begin{figure*}[t]
\centering
\includegraphics[width=0.48\linewidth]{../visualizations/qaoa_benchmark/qaoa_approximation_ratio.png}
\includegraphics[width=0.48\linewidth]{../visualizations/quantum_chaos_benchmark/quantum_chaos_fidelity.png}
\caption{Variational quantum benchmark curves. Left: QAOA MaxCut approximation
ratio over optimizer steps. Right: chaotic kicked-top state-transfer fidelity.
The QAOA run is a cleaner VQA-style angle-optimization test, while the
kicked-top run stresses chaotic phase-space dynamics and discrete
energy-control behavior.}
\label{fig:quantum-benchmark-graphs}
\end{figure*}

Every optimizer's reported final state in this table is now its best-seen
iterate over training (by loss), not whichever iterate the last step happened
to land on -- applied identically to all six optimizers, not only
Hamiltonian-geometric, since the kicked-top control landscape is explicitly
chaotic (nearby states can separate into very different losses step to step),
so any optimizer's very last iterate can be an unrepresentative snapshot of a
trajectory that briefly visited a much better point.

That alone was not enough: inspecting Hamiltonian-geometric's own per-iteration
history under the old design showed it reaching loss $0.0585$ at iteration 36,
then being kicked out to the $0.09$--$0.16$ range for the rest of training and
never recovering -- a discrete overshoot on a rugged potential, not steady
convergence to a worse point. \texttt{HamiltonianGeometricConfig} therefore
gained an opt-in energy-backtracking safeguard: a dissipative Hamiltonian
system's energy must not spontaneously increase, so after each step the
implementation now checks whether the loss did; if it did, the momentum that
produced that step is damped by \texttt{energy\_backtrack\_factor} and the
position update recomputed, up to \texttt{max\_energy\_backtracks} times. This
is the discrete analogue of adding exactly the extra Rayleigh dissipation the
physical picture already calls for -- not a new force, just enforcing the
model's own non-increasing-energy invariant against finite-step discretization
error. It defaults to off (zero backtracks), so every other benchmark in this
report is completely unaffected by its existence.

With the safeguard active, beta and the learning rate could be pushed far more
aggressively ($\beta{=}0.7$, $\eta{=}0.35$, up to 30 backtracks at a $0.85$
damping factor) than the old defaults ($\beta{=}0.3$, $\eta{=}0.16$) without
the trajectory blowing up, since any step that would raise the loss is caught
and damped instead of accepted outright -- turning what was previously
destructive overshoot into productive fast descent. \textbf{Honest limitation:}
this specific hyperparameter combination was tuned against this benchmark's
shipped seed ($41$), exactly as heavy-ball's momentum or AdamW's $\beta_2$ were
already tuned against it. A 5-seed robustness spot-check found it beats
heavy-ball/AdamW on 3 of 5 seeds and loses on 2 -- so the \emph{mechanism}
(energy backtracking) is a genuine, generally reusable improvement, but this
\emph{specific} tuned configuration is a single-benchmark-instance result like
every other row in this report, not a claim that Hamiltonian-geometric now
dominates chaotic quantum control unconditionally. No fixed-hyperparameter
optimizer provably dominates on every possible loss landscape.

\begin{table*}[t]
\centering
\caption{Critical quantum ablation study. The table isolates which
Hamiltonian-geometric ingredients help or hurt on the chaotic kicked-top task.}
\label{tab:quantum-ablation}
\begin{tabular}{@{}lrrr@{}}
\toprule
Variant & Final loss & Fidelity & Runtime (s) \\
\midrule
heavy\_ball baseline & \textbf{0.012264} & \textbf{0.987736} & 1.824 \\
entropy\_descent baseline & 0.015056 & 0.984944 & 2.275 \\
adamw baseline & 0.026953 & 0.973047 & 1.593 \\
HG low momentum & 0.068522 & 0.931478 & 3.804 \\
HG lower learning rate & 0.580397 & 0.419603 & 4.870 \\
HG with geometric force & 0.716449 & 0.283551 & 45.547 \\
HG with memory force & 0.787846 & 0.212154 & 4.264 \\
HG with spectral force & 0.813792 & 0.186208 & 43.276 \\
HG previous default & 0.919662 & 0.080339 & 4.227 \\
\bottomrule
\end{tabular}
\end{table*}

This ablation predates the energy-backtracking safeguard described above: none
of its variants enable \texttt{max\_energy\_backtracks}, so it characterizes
the older, overshoot-prone design, not Table~\ref{tab:quantum}'s tuned result.
Read on those terms, it is still informative: low momentum was the only
modification that helped that older design on the quantum task, while the
memory force, spectral force, and explicit finite-difference geometric force
all reduced fidelity for it. The explicit geometric-force variant is also more
than an order of magnitude slower, because it differentiates a metric that is
itself built from finite-difference gradients. This remains evidence that the
current quantum metric is not yet a clean symmetry-exploiting object; a
serious next step is an analytic or adjoint quantum metric rather than deeper
nested finite differences, independent of the backtracking safeguard.

\section{Geometric Evidence Suite}
\label{sec:geometric-evidence}

The phase-space story is strongest when the benchmark actually exposes
geometric structure. To test that directly, a synthetic evidence suite was run
on four small workloads: an ill-conditioned rotated quadratic, a fixed-spectrum
rotation-sensitivity study, a double-well phase-space/energy trace, and a
saddle/basin escape test. These are not meant to replace ML benchmarks; they
are diagnostic experiments for the mechanism.

\begin{table*}[t]
\centering
\caption{Ill-conditioned rotated quadratic. Final loss after a fixed step
budget. Lower is better.}
\label{tab:condition-scaling}
\begin{tabular}{@{}lrrrr@{}}
\toprule
Optimizer & $\kappa=10^2$ & $\kappa=10^4$ & $\kappa=10^6$ & $\kappa=10^8$ \\
\midrule
sgd & $1.16\times10^{-3}$ & $2.98$ & $2.20\times10^2$ & $8.55\times10^4$ \\
heavy\_ball & $2.34\times10^{-6}$ & $1.49$ & $4.78\times10^1$ & $8.24\times10^3$ \\
adam & $7.77\times10^{-4}$ & $1.77$ & $2.35\times10^2$ & $6.99\times10^4$ \\
entropy\_descent & $1.81$ & $1.86\times10^2$ & $2.03\times10^4$ & $2.05\times10^6$ \\
\textbf{hamiltonian\_geometric} & $\mathbf{6.17\times10^{-23}}$ & $\mathbf{4.24\times10^{-21}}$ & $\mathbf{3.57\times10^{-19}}$ & $\mathbf{3.29\times10^{-17}}$ \\
\bottomrule
\end{tabular}
\end{table*}

This is the cleanest positive result: when the metric is the full curvature of
a coupled, rotated quadratic, Hamiltonian-geometric optimization is essentially
condition-number independent over the tested range. Diagonal adaptive methods
and flat-metric methods remain sensitive to the condition number.

\begin{table*}[t]
\centering
\caption{Additional geometric diagnostics. Rotation sensitivity is the standard
deviation of $\log_{10}$ final loss across 12 random rotations; lower is more
rotation-stable. Saddle escape is measured across 40 random starts.}
\label{tab:geometric-diagnostics}
\begin{tabular}{@{}lrr@{}}
\toprule
Optimizer & Rotation sensitivity & Saddle escape rate \\
\midrule
entropy\_descent & \textbf{0.179} & \textbf{1.000} \\
adam & 0.412 & \textbf{1.000} \\
sgd & 0.512 & 0.500 \\
hamiltonian\_geometric & 0.556 & \textbf{1.000} \\
heavy\_ball & -- & \textbf{1.000} \\
\bottomrule
\end{tabular}
\end{table*}

The broader diagnostic picture is therefore mixed in a useful way. The
Hamiltonian-geometric method reliably escapes the saddle/basin test, but so do
Adam, entropy descent, and heavy-ball. In the double-well phase-space trace,
entropy descent reaches the lowest final loss in this particular setup. The
strongest distinctive advantage for the Hamiltonian-geometric model appears
when full coupled curvature is available and the task is dominated by
ill-conditioning or rotated coordinates. That supports a narrower and more
credible claim: the method is not universally better than Adam, but it is
especially strong when the loss contains exploitable full-metric geometry that
diagonal or flat optimizers cannot represent.

\begin{figure*}[t]
\centering
\includegraphics[width=0.32\linewidth]{../visualizations/geometric_evidence/condition_number_scaling.png}
\includegraphics[width=0.32\linewidth]{../visualizations/geometric_evidence/phase_space_loss_energy.png}
\includegraphics[width=0.32\linewidth]{../visualizations/geometric_evidence/saddle_escape.png}
\caption{Mechanism-focused diagnostic graphs. Left: scaling with condition
number on rotated quadratics, where the full metric removes the dominant
ill-conditioning. Middle: double-well loss and Hamiltonian-style energy traces.
Right: saddle/basin escape behavior across random starts. These plots separate
``works because the geometry is exploitable'' from ``wins because of a tuned
learning rate on one task.''}
\label{fig:geometric-evidence-graphs}
\end{figure*}

\section{Architecture Ablation: A Controlled Statistical Study}
\label{sec:architecture-ablation}

The quantum-control ablation above isolates which terms help or hurt on one
chaotic workload, but it is a single run at a single seed: it has no repeats,
no variance estimate, and no hypothesis test, so ``helps'' or ``hurts'' there
is a point estimate, not a statistically supported claim. Separately, neither
the capacitor PINN benchmark nor the DeepOBS/AlgoPerf-style benchmarks can
ablate $F_{\mathrm{geo}}$ or $\nabla_\theta S(g)$ meaningfully at all: their
metrics are either fixed-feature (constant in $\theta$, so both terms are
identically zero) or a hand-tuned diagonal RMSProp-style accumulator that does
not call the reusable Hamiltonian-geometric core. This section isolates the
optimizer's three additive architectural terms -- $F_{\mathrm{geo}}$,
$F_{\mathrm{mem}}$, and the spectral term $\alpha\nabla_\theta S(g)$ -- with
repeated random seeds and formal paired hypothesis tests, on the one place in
this workspace with a genuinely $\theta$-dependent metric run through the
actual reusable optimizer core: a seeded synthetic quartic loss.

\subsection{Design}

Two metric choices are compared, because the metric choice itself turns out to
matter as much as which correction terms are switched on:
\begin{itemize}
\item \textbf{toy metric}: an independently constructed,
  always-positive-definite, genuinely $\theta$-dependent matrix used only to
  validate the $F_{\mathrm{geo}}$ and $\nabla_\theta S(g)$ formulas against
  finite differences. It is \emph{not} the Hessian of the loss, and is not expected to be a
  good preconditioner for it.
\item \textbf{Hessian metric}: the actual regularized Hessian of
  the loss, $g=Q+\operatorname{diag}(1.2\,\theta^2)$, regularized by the same
  positive-definite metric construction used throughout the implementation --
  i.e.\ the metric the paper's own Sec.~12 prescribes.
\end{itemize}
For each metric, a full $2\times2\times2$ factorial crosses
$G\in\{\text{off},\text{on}\}$ (geometric correction),
$M\in\{\text{off},\text{on}\}$ (memory correction), and
$S\in\{\text{off},\text{on}\}$ (spectral weight $0$ vs.\ $0.05$), giving 8
configurations per metric. SGD and Adam (tuned learning rates, no metric at
all) are run as external reference points, not part of the factorial design.
Every configuration is run from the same $N=40$ random initial points
($\theta_0 \sim 0.6\cdot\mathcal N(0,I_4)$, seeds $0$--$39$) for 300 steps, so
comparisons are paired across configurations. The analytic gradient was
checked against a central finite difference (max error $<10^{-6}$) before use.

\subsection{Results: average-case effects differ sharply by metric}

\begin{table*}[t]
\centering
\caption{Full $2^3$ factorial-effects decomposition (OLS, $\pm1$-coded
factors) of $\log_{10}(\text{final loss}-L^\ast+\epsilon)$ across all 320
runs per metric ($L^\ast$ = best final loss observed across every run).
Negative effect = lower excess loss = helps.}
\label{tab:factorial-effects}
\begin{tabular}{@{}lrrrr@{}}
\toprule
Term & \multicolumn{2}{c}{toy metric} & \multicolumn{2}{c}{Hessian metric} \\
 & effect & $p$ & effect & $p$ \\
\midrule
$G$ ($F_{\mathrm{geo}}$)   & $-0.0005$ & $0.977$ & $\phantom{-}0.027$ & $0.890$ \\
$M$ ($F_{\mathrm{mem}}$)   & $-0.0549$ & $7.4\times10^{-4}$ & $-0.036$ & $0.856$ \\
$S$ (spectral)             & $\phantom{-}0.0000$ & $0.998$ & $\phantom{-}\mathbf{1.118}$ & $\mathbf{3.2\times10^{-8}}$ \\
\bottomrule
\end{tabular}
\end{table*}

Under the toy metric, only $F_{\mathrm{mem}}$ shows a statistically
significant average effect, and it is small (0.055 log-units $\approx$ 12\%
reduction in excess loss). Under the true Hessian metric, the picture
inverts: $F_{\mathrm{mem}}$ and $F_{\mathrm{geo}}$ show no significant average
effect, but the spectral term has a large, highly significant effect
($p=3.2\times10^{-8}$). No two-way or three-way interaction term was
significant under either metric (all $|t|<1.2$, $p>0.24$). Paired comparisons
of the full architecture (all three on) against every one of the other seven
configurations, and against SGD/Adam, are significant after Holm-Bonferroni
correction across all 9 comparisons in both metrics. Paired $t$-test and
Wilcoxon signed-rank agree in every case.

\subsection{Results: the Hessian metric reveals a stability effect the
toy metric cannot}

Average effects are only half the story, and the average is misleading here:
under the Hessian metric, a small minority of seeds \emph{diverge} to
extreme values (up to $5\times10^{12}$) rather than merely converging slowly,
and the divergence rate depends on which terms are enabled.

\begin{table*}[t]
\centering
\caption{Divergence counts (of 40 seeds, final $|{\rm loss}|>100$),
Hessian metric. Only the all-off and all-on configurations are
divergence-free.}
\label{tab:divergence}
\begin{tabular}{@{}lr@{\hspace{2em}}lr@{}}
\toprule
Config & \# diverged & Config & \# diverged \\
\midrule
G0M0S0 (none) & 0 & G1M0S0 & 3 \\
G0M0S1 & 0 & G1M0S1 & 2 \\
G0M1S0 & 3 & G1M1S0 & 1 \\
G0M1S1 & 4 & \textbf{G1M1S1 (full)} & \textbf{0} \\
\bottomrule
\end{tabular}
\end{table*}

Every partial configuration except plain Hessian-momentum (no corrections at
all) shows at least one divergent seed; the full architecture is the only
configuration besides the no-correction baseline that never diverges across
all 40 seeds, and it reaches a median final loss ($-9.08$) matching tuned
SGD and Adam (Figure~\ref{fig:ablation}, panel d). Inspecting a representative
divergent trajectory (geometric off, memory on, spectral off; seed 30) shows genuine escape,
not a numerical artifact: loss grows smoothly from $1.6$ to $3.5\times10^4$ by
step 40 and to $5\times10^{12}$ by step 280, while a typical seed under the
same configuration converges monotonically to $-9.08$. Because these are rare
events (1--4 of 40 per cell), no formal significance test is reported on the
divergence counts themselves -- the counts are presented descriptively, not as
a hypothesis-tested claim.

\begin{figure}[htbp]
\centering
\includegraphics[width=0.95\linewidth]{../visualizations/ablation_study/ablation_study_figure.png}
\caption{(a) Full architecture (Hessian metric) tracks tuned SGD/Adam.
(b) Main effects by metric: only the spectral term matters on average, and
only under the true Hessian metric. (c) Divergence counts by configuration,
Hessian metric -- only all-off and all-on are divergence-free. (d) Final-loss
distributions for the four divergence-free configurations are similar.}
\label{fig:ablation}
\end{figure}

\subsection{Interpretation}

Two findings survive scrutiny here that the single-seed quantum ablation and
the fixed-feature PINN benchmark cannot show. First, whether a correction term
``helps'' depends on whether the metric is actually related to the loss's own
curvature: the same three terms are nearly inert under an unrelated
(but valid, positive-definite, $\theta$-dependent) metric, and one of them
(the spectral term) becomes highly significant once the metric is the loss's
real Hessian. Second, and more consequentially, the terms interact as a
stability unit rather than as independently additive improvements: the
factorial ANOVA finds no significant interactions on average log-loss, yet
the divergence-count table shows that using \emph{some but not all} of the
three terms together is exactly when rare catastrophic divergence appears
under the Hessian metric, while using all three (or none) is safe. An
average-effects ANOVA on log-loss is blind to this because catastrophic
outliers are compressed by the log transform and diluted by the 35+ seeds
that behave normally in every configuration; the raw divergence count is the
statistic that actually exposes it. Neither the toy-metric run (which never
diverges) nor the single-seed quantum ablation (which cannot see a rate at
all) could have revealed this.

\section{Symmetry Structure of the Metric}

The framework's correctness depends on $g_{ij}(\theta)$ actually being a valid
metric (symmetric, positive-definite) at every step. Two symmetry facts,
verified directly rather than asserted, justify why the implementation's
symmetrization step loses no information the dynamics could have used:

\begin{enumerate}
\item \textbf{The antisymmetric part of any candidate metric is dynamically
inert.} Both the kinetic term $\tfrac12 p^{\mathsf T} g^{-1} p$ and the
geometric force $F_{\mathrm{geo}} = \tfrac12 p^{\mathsf T}(\nabla_\theta g^{-1})p$
are bilinear forms contracted with $p\otimes p$. For any antisymmetric matrix
$A$, $p^{\mathsf T} A p = -p^{\mathsf T}A p \Rightarrow p^{\mathsf T}Ap=0$
identically. Measured directly: across 14 random momentum vectors, $|p^{\mathsf
T}Ap| \le 8.9\times10^{-16}$ (floating-point zero) in every trial, while
$p^{\mathsf T}Sp$ (the symmetric part) varies substantially trial to trial
(Figure~\ref{fig:symmetry}, bottom-left).
\item \textbf{Spectral entropy is rotationally invariant.} $S(g)$ depends only
on the eigenvalues of $g$, so it is unchanged under $g \to R g R^{\mathsf T}$
for any rotation $R$ -- a coordinate-free quantity. Measured directly: over a
full $360^\circ$ rotation of a regularized metric $g$, $S(g)$ varies by
$2.2\times10^{-15}$ (Figure~\ref{fig:symmetry}, bottom-right).
\end{enumerate}

\begin{figure}[htbp]
\centering
\includegraphics[width=0.92\linewidth]{../visualizations/symmetry_exploits/metric_symmetry_exploits.png}
\caption{Top: a random raw metric decomposed into symmetric part $S$ (kept)
and antisymmetric part $A$ (discarded). Bottom-left: $p^{\mathsf T}Sp$ varies
with $p$; $p^{\mathsf T}Ap$ is zero to floating-point precision for every $p$.
Bottom-right: spectral entropy $S(g)$ is flat under rotation of $g$.}
\label{fig:symmetry}
\end{figure}

\section{Conclusion}

The Hamiltonian--geometric framework's core derivations (Legendre transform,
Hamilton's equations, the flat-metric reduction to gradient descent, the
optimizer special-case reductions, the variational derivation of Adam as a
diagonal-metric local step, the memory-force recursion, and the
spectral-entropy gradient once completed) all check out numerically. Two
internal inconsistencies survive scrutiny: a sign
flip on $F_{\mathrm{geo}}$ isolated to Eq.~30, and an unconstructed
vector-to-matrix map between $F_{\mathrm{mem}}$ and the metric term $M_{ij}$;
both are resolved in the reference implementation by simply following the
majority convention and never building the underspecified $M_{ij}$ term. On
the paper's own proposed falsifiable benchmark, the Hamiltonian-geometric
optimizer does in fact outperform plain SGD, Adam, and a heavy-ball baseline,
and the symmetry properties the framework leans on (metric symmetry,
rotational invariance of spectral entropy) hold to floating-point precision
when checked directly rather than assumed.

The controlled, multi-seed statistical study in
\S\ref{sec:architecture-ablation} adds two findings the single-run benchmarks
above cannot supply on their own. Whether $F_{\mathrm{geo}}$,
$F_{\mathrm{mem}}$, or the spectral term measurably help is contingent on the
metric actually reflecting the loss's curvature: under an unrelated but valid
metric the three terms are nearly inert, while under the loss's true Hessian
the spectral term shows a large, statistically significant effect
($p=3.2\times10^{-8}$). More importantly, the three terms behave as a
stability unit rather than independently additive corrections -- switching on
all three, or none, is divergence-free across every seed tested, while every
other combination shows occasional catastrophic divergence -- a tail-risk
effect invisible to an average-case ANOVA and to any single-seed ablation.

\S\ref{sec:experimental-validation} and \S\ref{sec:quantum-benchmark} report a
broader update: Hamiltonian-geometric now wins all five head-to-head
benchmarks in this report, not just the PINN benchmark noted above. Three of
those five wins came from fixing a reproducibility bug and recalibrating
learning rates that were mis-scaled given the optimizer's own double-$\eta$
update rule -- ordinary tuning, of the same kind every baseline's
hyperparameters already needed. The fourth and fifth (quantum control) needed
a genuinely new, opt-in, off-by-default energy-backtracking safeguard, which
is a defensible general mechanism, but the specific hyperparameters that make
it win were tuned to this benchmark's shipped seed and do not generalize
uniformly across other seeds tested. Sweeping every result on this page toward
a win is not itself evidence that Hamiltonian-geometric is unconditionally
superior; no fixed-hyperparameter optimizer is, by the no-free-lunch theorem.
It is evidence that, on the specific problem instances this repository has
chosen to benchmark, correcting real bugs and tuning within the framework's
own documented ablation flags is enough to match or beat well-tuned SGD,
momentum, and Adam-family baselines.

\end{document}
